\section{Carta PC--Fock completa, memoria y almacenamiento terminal}
\label{sec:ns-pcf-prospective-construction}

Fijamos \(s>5/2\), \(\nu>0\), un corte de Galerkin \(M\), una profundidad
nonádica \(J\) y un horizonte \(T<\infty\). Sea
\(V_M\subset H^s_\sigma(\mathbb T^3)\) el espacio solenoidal de media cero y
sean
\begin{equation}
 A_M=P_M(-\Delta)P_M,
 \qquad
 N_M(u)=P_M\mathbb P(u\cdot\nabla u).
 \label{eq:ns-import-galerkin-operators}
\end{equation}
Esta sección construye el portador PC--Fock completo, sus lectores, memorias,
acarreo, almacenamiento terminal y retorno, que ingresan en la realización
propietaria de la sección siguiente. Su KYP local, su observabilidad NS--OBS y
las dominaciones LRDN--LCN son exclusivamente
\texttt{CONTROL\_NO\_GOBERNANTE}: no condicionan la métrica
\(U^{\rm own}\), su telescopía, su presupuesto raíz ni la conclusión global.

\subsection{Torre completa y lectores de grados uno y dos}

La elevación física se toma en la Fock simétrica completa:
\begin{equation}
 \mathcal V_M^{\rm PCF}(u)
 :=\bigoplus_{n\ge0}\frac{R_{\rm F}^n}{\sqrt{n!}}u^{\odot n},
 \qquad
 \|\mathcal V_M^{\rm PCF}(u)\|^2
 =e^{R_{\rm F}^2\|u\|_{H^s}^2}.
 \label{eq:ns-import-pcf-lift}
\end{equation}
La torre no se trunca. Para \(z_{n,M}=u_M^{\odot n}\), la regla de Leibniz
da la recurrencia exacta
\begin{equation}
 \dot z_{n,M}
 =L_{n,M}z_{n,M}
 +\mathcal N_{n,n+1;M}z_{n+1,M}
 +\mathcal F_{n,n-1;M}(f_M)z_{n-1,M},
 \qquad n\ge1,
 \label{eq:ns-import-carleman-all-degrees}
\end{equation}
donde, sobre la diagonal coherente,
\begin{align}
 L_{n,M}u^{\odot n}
 &=-\nu n(A_Mu)\odot u^{\odot(n-1)}, \notag\\
 \mathcal N_{n,n+1;M}u^{\odot(n+1)}
 &=-nN_M(u)\odot u^{\odot(n-1)}, \notag\\
 \mathcal F_{n,n-1;M}(f)u^{\odot(n-1)}
 &=nf\odot u^{\odot(n-1)}.
 \label{eq:ns-import-carleman-blocks}
\end{align}

El espacio de estado es la suma ortogonal
\begin{equation}
 \mathscr K_{M,j}^{\rm PCF}
 =\mathscr K_{M,j}^{\rm HMT}\widehat\otimes\Gamma_s(V_M)
 \oplus\mathscr K_{M,j}^{\kappa}
 \oplus\mathscr K_{M,j}^{\rm micro}
 \oplus\mathscr K_{M,j}^{\rm sh},
 \label{eq:ns-import-full-carrier}
\end{equation}
donde \(\mathscr K_{M,j}^{\rm HMT}\) conserva hoja, orientación,
\(q_A,q_V,T\), borde, ruta y las dos memorias \(D/H\). Los grados uno y dos
se leen mediante
\begin{equation}
 \begin{aligned}
 z_M(u)&=\nu^{1/2}A_M^{1/2}\Lambda^su,\\
 w_M(u)&=\nu^{-1/2}A_M^{-1/2}\Lambda^sN_M(u),\\
 \mathcal Y_{\beta,M}\mathcal V_M^{\rm PCF}(u)
 &=\sqrt\beta\,z_M(u)
   \oplus\frac{1}{2\sqrt\beta}\,w_M(u).
 \end{aligned}
 \label{eq:ns-import-degree-one-two-reader}
\end{equation}
En consecuencia,
\begin{equation}
 |B_{s,M}(u)|
 \le
 \|\mathcal Y_{\beta,M}\mathcal V_M^{\rm PCF}(u)\|^2
 =\beta\nu D_{s,M}(u)+\frac{\|w_M(u)\|^2}{4\beta}.
 \label{eq:ns-import-young-port}
\end{equation}
La segunda coordenada es un operador lineal uniformemente acotado sobre
\(\operatorname{Sym}^2H^s_\sigma\):
\begin{equation}
 \sup_M\|\mathfrak C_M\|
 \le\widetilde M_{s,\beta,\nu},
 \qquad
 C_M^{\rm obs}:=\sqrt2\,\mathfrak C_M\pi_2,
 \qquad
 \Sigma_M:=\bigl(I+(C_M^{\rm obs})^*C_M^{\rm obs}\bigr)^{-1}
 \succeq\frac{I}{1+2\widetilde M_{s,\beta,\nu}^2}.
 \label{eq:ns-import-uniform-graph}
\end{equation}

\subsection{Memoria bidireccional y almacenamiento terminal}

Con
\begin{equation}
 r=\frac{\pi_{\rm HMT}}{\varphi_{\rm HMT}^2},
 \qquad
 D_{M,j}^{\rm mem}=(r-1)(M_{M,j}^+-M_{M,j}^-),
 \qquad
 H_{M,j}=(1+r)(M_{M,j}^++M_{M,j}^-),
 \label{eq:ns-import-dh}
\end{equation}
las dos memorias satisfacen
\begin{equation}
 \dot D_{M,j}^{\rm mem}=B_{s,M,j},
 \qquad
 \dot H_{M,j}=\frac{1+r}{r-1}|B_{s,M,j}|.
 \label{eq:ns-import-dh-dot}
\end{equation}
Poniendo
\[
 \delta_r=\frac{r-1}{r+1},
 \qquad
 \mathcal N_{M,j}:=\delta_rH_{M,j}-D_{M,j}^{\rm mem}
 =2(r-1)M_{M,j}^-\ge0,
\]
se obtiene
\begin{equation}
 \dot{\mathcal N}_{M,j}=|B_{s,M,j}|-B_{s,M,j}.
 \label{eq:ns-import-negative-memory}
\end{equation}
El almacenamiento terminal explícito es
\begin{equation}
 \begin{aligned}
 G_{M,j,T}(t)
 &=E_{s,M,j}(t)
   +2(r-1)\bigl(M_{M,j}^-(T)-M_{M,j}^-(t)\bigr),\\
 G_{M,j,T}(t)&\ge E_{s,M,j}(t),\\
 \dot G_{M,j,T}+\nu D_{s,M,j}+|B_{s,M,j}|
 &=F_{s,M,j}.
 \end{aligned}
 \label{eq:ns-import-terminal-storage}
\end{equation}
Si
\[
 h^-_{M,j;t,T}(\tau)
 :=\mathbf1_{[t,T]}(\tau)\sqrt{\dot M^-_{M,j}(\tau)},
\]
la columna
\begin{equation}
 \mathcal E_{M,j,T}^{\rm term}(t)\Xi
 :=2^{-1/2}\Lambda^su_{M,j}(t)
 \oplus\sqrt{2(r-1)}\,h^-_{M,j;t,T}
 \label{eq:ns-import-terminal-column}
\end{equation}
realiza el almacenamiento como un Gram:
\begin{equation}
 \|\mathcal E_{M,j,T}^{\rm term}(t)\Xi\|^2
 =G_{M,j,T}(t).
 \label{eq:ns-import-terminal-gram}
\end{equation}

El acarreo común entre profundidades es
\begin{equation}
 \kappa_{M,j}
 =\frac12\left(E_j|B_{s,M,j+1}|-|B_{s,M,j}|\right)\ge0,
 \qquad
 E_jB_{s,M,j+1}=B_{s,M,j}.
 \label{eq:ns-import-carry}
\end{equation}
Su coordenada de salida se fija por
\begin{equation}
 q_{M,j}^{\kappa}
 :=\left(\frac{2(1+r)}{r-1}\kappa_{M,j}\right)^{1/2},
 \qquad
 \|q_{M,j}^{\kappa}\|^2
 =\frac{2(1+r)}{r-1}\kappa_{M,j}.
 \label{eq:ns-import-carry-port}
\end{equation}
Al mismo tiempo,
\begin{equation}
 E_j\Delta_T\mathcal N_{M,j+1}
 =\Delta_T\mathcal N_{M,j}+2K_{M,j}(T),
 \qquad
 K_{M,j}(T):=\int_0^T\kappa_{M,j}(t)\,dt.
 \label{eq:ns-import-carry-telescope}
\end{equation}
Las innovaciones microscópicas y las bandas de Galerkin permanecen
ortogonales:
\begin{equation}
 Q_{\rm micro}^{\rm APP}=I-J_9E_9,
 \qquad E_9y^{\rm micro}=0,
 \qquad \mathsf K_{\rm carry}y^{\rm micro}=0,
 \label{eq:ns-import-micro}
\end{equation}
y, grado a grado,
\begin{equation}
 \mathcal N_{n,n+1;3M}J_M^{(n+1)}
 =J_M^{(n)}\mathcal N_{n,n+1;M}
 +\mathcal S_{n,M}^{\rm sh},
 \qquad
 E_M^{(n)}\mathcal S_{n,M}^{\rm sh}=0.
 \label{eq:ns-import-shell}
\end{equation}

\subsection{Control KYP finito no gobernante y cascada de la carta}

La fuerza y su término directo se incorporan mediante
\begin{equation}
 \begin{aligned}
 x_{M,j+1}
 &=\mathcal A_{M,j}^{[j]}x_{M,j}
   +\mathcal B_{M,j,T}^{F,[j]}g_{M,j},\\
 y_{M,j}^{\rm loss}
 &=\mathcal C_{M,j,T}^{[j]}x_{M,j}
   +\mathcal D_{M,j,T}^{F,[j]}g_{M,j},
 \end{aligned}
 \label{eq:ns-import-affine-step}
\end{equation}
con la descomposición ortogonal
\begin{equation}
 \mathcal C_{M,j,T}
 =\mathcal C_{M,j,T}^{(1)}
 \oplus\mathcal C_{M,j,T}^{(2)}
 \oplus\mathcal C_{M,j,T}^{\kappa}
 \oplus\mathcal C_{M,j,T}^{\rm micro}
 \oplus\mathcal C_{M,j,T}^{\rm sh}.
 \label{eq:ns-import-loss-decomposition}
\end{equation}
El término directo aparece ya en
\[
 g_M=A_M^{-1/2}\Lambda^sP_Mf,
 \qquad
 a_M=\frac{g_M}{2\sqrt\nu},
 \qquad
 y_{D,M}=\sqrt\nu A_M^{1/2}\Lambda^su_M-a_M,
\]
para el que
\begin{equation}
 dG_{M,j,T}
 +\bigl(\|y_{D,M}\|^2+|B_{s,M,j}|\bigr)\,dt
 =\|a_M\|^2\,dt.
 \label{eq:ns-import-force-feedthrough}
\end{equation}

La afinidad se homogeneiza sin duplicar el dato. Sea
\[
 \mathscr G_M^{F,[j]}
 :=\bigoplus_{k=j}^{J-1}L^2(\Omega_k;\mathscr G^F_{M,k,T}),
 \qquad
 \mathbf x_{M,j}:=x_{M,j}\oplus(g_{M,j},\ldots,g_{M,J-1}),
\]
y sean \(e_j\) la evaluación de la primera entrada y \(\mathsf S_j\) el
desplazamiento de la cola. Definimos
\begin{equation}
 \mathbf\Gamma_{M,j}
 :=\begin{pmatrix}
  \mathcal A_{M,j}^{[j]}&\mathcal B_{M,j,T}^{F,[j]}e_j\\
  0&\mathsf S_j
 \end{pmatrix},
 \qquad
 \mathbf L_{M,j,T}
 :=\mathcal Q_{M,j,T}^{\rm loss,inc}
 \begin{pmatrix}
  \mathcal C_{M,j,T}^{[j]}&\mathcal D_{M,j,T}^{F,[j]}e_j
 \end{pmatrix}.
 \label{eq:ns-import-homogeneous-affine-step}
\end{equation}
Aquí \(\mathcal Q_{M,0,T}^{\rm loss,inc}=I\); para \(j>0\) elimina
únicamente la copia heredada de los canales PC--Fock y actúa como la identidad
sobre acarreo, innovación microscópica y banda espectral.

\begin{theorem}[Control de factorización KYP finita no circular]
\label{thm:ns-import-affine-stein}
En el nivel terminal se define
\begin{equation}
 \mathbf U_{M,J,T}(t)
 :=(\mathcal E_{M,J,T}^{\rm term}(t))^*
   \mathcal E_{M,J,T}^{\rm term}(t).
 \label{eq:ns-import-terminal-metric}
\end{equation}
Fijados \(M,J,T\), sean \(A,B,C,D\) los cuatro bloques del paso afín y
\(U_+=\widetilde{\mathbf U}_{M,j+1,T}(t)\). Defínanse
\begin{equation}
 Q:=I-B^*U_+B-D^*D,
 \qquad L:=B^*U_+A+D^*C.
 \label{eq:ns-import-kyp-q-l}
\end{equation}
Si \(Q\succ0\), la recurrencia no circular es
\begin{equation}
 U:=A^*U_+A+C^*C+L^*Q^{-1}L.
 \label{eq:ns-import-affine-stein}
\end{equation}
Entonces la matriz KYP completa satisface
\begin{equation}
 \begin{aligned}
 \mathfrak S_{M,j,T}
 &:=
 \begin{pmatrix}
 U-A^*U_+A-C^*C&-A^*U_+B-C^*D\\
 -B^*U_+A-D^*C&I-B^*U_+B-D^*D
 \end{pmatrix}\\
 &=
 \begin{pmatrix}Q^{-1/2}L&-Q^{1/2}\end{pmatrix}^{*}
 \begin{pmatrix}Q^{-1/2}L&-Q^{1/2}\end{pmatrix}
 \succeq0.
 \end{aligned}
 \label{eq:ns-import-affine-supply}
\end{equation}
Para \(Q\succeq0\) vale la variante con pseudoinversa si
\(\operatorname{Ran}L\subseteq\operatorname{Ran}Q^{1/2}\). Esta construcción
es finita: no afirma que \(Q^{-1}\), \(U\) o la ganancia permanezcan
uniformes al retirar \(M,J\).
\end{theorem}

\begin{proof}
Por \eqref{eq:ns-import-affine-stein}, el bloque superior izquierdo es
\(L^*Q^{-1}L\); los bloques no diagonales son \(-L^*\) y \(-L\), y el bloque
inferior derecho es \(Q\). La multiplicación de la fila mostrada en
\eqref{eq:ns-import-affine-supply} da exactamente los cuatro bloques. La
positividad precede así a la toma de cualquier raíz.
\end{proof}

Este teorema comprueba únicamente la carta finita cuando se satisfacen sus
hipótesis propias. No aporta una premisa al owner Stein--Lyapunov ni limita la
uniformidad demostrada por su Gram terminal y su raíz común.

En las coordenadas homogeneizadas se toma
\(\mathbf U_{M,j,T}=U\) y se denota por
\(\mathbf R_{M,j,T}\) el levantamiento de la fila
\(\begin{psmallmatrix}Q^{-1/2}L&-Q^{1/2}\end{psmallmatrix}\).
Con esta convención se recupera, para cada truncación finita que satisface la
hipótesis KYP,
\[
 \mathbf U_{M,j,T}
 =\mathbf\Gamma_{M,j}^*\widetilde{\mathbf U}_{M,j+1,T}
  \mathbf\Gamma_{M,j}
 +\mathbf L_{M,j,T}^*\mathbf L_{M,j,T}
 +\mathbf R_{M,j,T}^*\mathbf R_{M,j,T}.
\]

La columna
\begin{equation}
 [\mathbf x]_{\mathbf U_{M,j,T}}
 \longmapsto
 [\mathbf\Gamma_{M,j}\mathbf x]_{
  \widetilde{\mathbf U}_{M,j+1,T}}
 \oplus\mathbf L_{M,j,T}\mathbf x
 \oplus\mathbf R_{M,j,T}\mathbf x
 \label{eq:ns-import-local-isometry}
\end{equation}
es isométrica. Su dilatación de Julia y su composición producen
\(\mathfrak U_{M,J,T}\) con
\begin{equation}
 \mathfrak U_{M,J,T}^{\sharp}\mathfrak U_{M,J,T}=I
 \label{eq:ns-import-cascade-isometry}
\end{equation}
y el telescopio
\begin{equation}
 \begin{aligned}
 \|\mathbf x_{M,0}\|_{\mathbf U_{M,0,T}}^2
 ={}&\mathbb E_{0:J}
 \|\mathbf x_{M,J}\|_{\mathbf U_{M,J,T}}^2\\
 &+\sum_{j=0}^{J-1}\mathbb E_{0:j}
 \left(\|\mathbf L_{M,j,T}\mathbf x_{M,j}\|^2
       +\|\mathbf R_{M,j,T}\mathbf x_{M,j}\|^2\right).
 \end{aligned}
 \label{eq:ns-import-stein-telescope}
\end{equation}

\subsection{Presupuesto de la carta y retorno de la velocidad}

Sobre la diagonal física, la primera coordenada terminal es
\(G_{M,J,T}\), y las demás coordenadas forman la suma de pérdidas:
\begin{equation}
 \|Z_{M,J,t}\|^2
 =\mathbb E_{0:J}G_{M,J,T}(t)
 +\sum_{j=0}^{J-1}\mathbb E_{0:j}\Lambda_{M,j}([0,t]).
 \label{eq:ns-import-ledger}
\end{equation}
Como \(q_{M,j}^{\kappa}\) es una de esas coordenadas,
\eqref{eq:ns-import-carry-telescope} se integra una sola vez:
\begin{equation}
 2\sum_{j=0}^{J-1}\mathbb E_{0:j}K_{M,j}(T)
 \le\|\Xi_{M,0}\|^2+Q_f(T).
 \label{eq:ns-import-carry-funded}
\end{equation}
De aquí,
\begin{equation}
 \mathbb E_{0:J}G_{M,J,T}(0)
 \le E_{s,M}(P_Mu_0)
 +\Delta_T\mathcal N_0+\|\Xi_{M,0}\|^2+Q_f(T).
 \label{eq:ns-import-root-bound}
\end{equation}
La raíz es común a todos los cortes:
\begin{equation}
 \Xi_{M,0}=J_{0\to M}\Xi_0,
 \qquad J_{0\to M}^*J_{0\to M}=I,
 \qquad \mathcal N_{M,0}=\mathcal N_0.
 \label{eq:ns-import-common-root}
\end{equation}
Además,
\begin{equation}
 \|\Xi_{M,0}\|^2
 \le C_{\rm HMT}
 +(1+2\widetilde M_{s,\beta,\nu}^2)
 e^{R_{\rm F}^2\|u_0\|_{H^s}^2},
 \qquad
 Q_f(T)\le C_{\nu,s}\int_0^T\|f(t)\|_{H^{s-1}}^2\,dt.
 \label{eq:ns-import-root-data-bound}
\end{equation}
Por tanto, para esta carta y estos \(M,J,T\), ponemos
\begin{align}
 \mathcal B_{M,J,T}^{\rm route}:={}&
 \frac12\|u_0\|_{H^s}^2
 +|\Delta_T\mathcal N_{M,0}|+C_{\rm HMT}
 +(1+2\widetilde M_{s,\beta,\nu}^2)
 e^{R_{\rm F}^2\|u_0\|_{H^s}^2}\notag\\
 &+C_{\nu,s}\int_0^T\|f(t)\|_{H^{s-1}}^2\,dt
 <\infty
 \label{eq:ns-import-explicit-budget}
\end{align}
La presencia explícita de \(\Delta_T\mathcal N_{M,0}\) impide usar esta
expresión como prueba independiente de uniformidad. Es un diagnóstico local
de la carta; no reemplaza el presupuesto owner
\eqref{eq:nsclose-owner-data-budget}. La identidad terminal y la desigualdad
de Young dan, para el corte fijo,
\begin{equation}
 \boxed{
 \sup_{0\le t\le T}\|Z_{M,J,t}\|^2
 +\frac{\nu}{2}\int_0^TD_{s,M}(t)\,dt
 \le\mathcal B_{M,J,T}^{\rm route}.}
 \label{eq:ns-import-uniform-budget}
\end{equation}

La publicación de la velocidad se normaliza por
\begin{equation}
 \iota_{s,M}u=2^{-1/2}\Lambda^su,
 \qquad
 \mathcal R_{1,M}=\sqrt2\,\Lambda^{-s}\pi_1,
 \qquad
 \mathcal R_{1,M}\iota_{s,M}=I_{V_M}.
 \label{eq:ns-import-grade-one-return}
\end{equation}
El lector de la cascada es
\begin{equation}
 \mathcal R_{M,J,T}^{\rm phys}
 :=\mathcal R_{1,M}\mathfrak U_{M,J,T}^{\sharp}.
 \label{eq:ns-import-full-return}
\end{equation}
Por \eqref{eq:ns-import-cascade-isometry},
\begin{equation}
 \boxed{
 \mathcal R_{M,J,T}^{\rm phys}
 \mathfrak U_{M,J,T}\mathcal V_M^{\rm PCF}(u_M(t))
 =u_M(t).}
 \label{eq:ns-import-exact-return}
\end{equation}
El retorno emplea el adjunto métrico de la misma cascada y recupera la
publicación hidrodinámica sin seleccionar una rama probabilística.
